\section{The two-stock benchmark and the demand-floor law}\label{sec:benchmark}

\paragraph{Model.} Two interacting stocks are harvested under an aggregate observation and a demand floor:
\begin{equation}\label{eq:bench}
\dot x_i=r_ix_i\Bigl(1-\frac{x_i}{10}\Bigr)-\frac{x_1x_2}{20}-u_i+d_i,\qquad i=1,2,
\end{equation}
with $r=(1,\tfr45)$, shocks $|d_i|\le\tfr1{10}$, constraints $q_i(x)=x_i-2\ge0$ (conservation floors), and the action set
\begin{equation}\label{eq:benchU}
U_D=\bigl\{u\in[0,2]^2:\ u_1+u_2\ge D\bigr\},\qquad D=2\ \text{unless stated otherwise},
\end{equation}
a polytope whose last facet is the demand floor. The controller observes only the aggregate stock $Y=x_1+x_2$, so its atom on the band $Y\ge4$ is the fibre $B_Y=\{x\in\calS:\ x_1+x_2=Y\}$. Along the fibre the constraint $q_i$ is active only at the corner $x^{(i)}$ where $x_i=2$; interior points of the fibre have no active constraint and contribute no row. The active bundle $\calK(B_Y)$ therefore consists of the rows of the two corners, one per shock value, and the worst shock $d_i=-\tfr1{10}$ gives the binding rows
\begin{equation}\label{eq:caps}
\psi_i(u)=\capi{i}(Y)-u_i,\qquad \capi{1}(Y)=\tfr32-\tfr{Y-2}{10},\qquad \capi{2}(Y)=\tfr{59}{50}-\tfr{Y-2}{10},
\end{equation}
obtained by evaluating \eqref{eq:bench} at $x^{(1)}=(2,Y-2)$ and $x^{(2)}=(Y-2,2)$: $\capi{1}(Y)=1\cdot2\cdot\tfr45-\tfr{2(Y-2)}{20}-\tfr1{10}$ and $\capi{2}(Y)=\tfr45\cdot2\cdot\tfr45-\tfr{2(Y-2)}{20}-\tfr1{10}$. Both caps lie in $[0,2]$ for $Y\in[4,\tfr{69}5]$, so the box facets of $U_D$ are never charged on that band and the only constraint of $U_D$ that can bind is the floor. The fibre is \emph{first-order viable under the aggregate observation} when $\calR(B_Y)=\{u\in U_D:\ u_1\le\capi{1}(Y),\ u_2\le\capi{2}(Y)\}\neq\emptyset$. (The first-order test is necessary for viability; it is not claimed to be sufficient along the whole band.)

\begin{proposition}[fibre margin, certificate and Farkas pair]\label{prop:benchmark}
For $D=2$ and $Y\in[4,\tfr{69}5]$:
\begin{enumerate}[leftmargin=1.6em,itemsep=1pt]
\item[(i)] the lower value of the fibre is $\varepsilon^\ast(Y)=\max_{u\in U_2}\min_i\psi_i(u)=\tfr{27-5Y}{50}$, attained at the $Y$-independent action $u^\ast=(\tfr{29}{25},\tfr{21}{25})$; the fibre is first-order viable iff $Y\le Y^\ast=\tfr{27}5$, equivalently iff its width $Y-4$ is at most $\tfr75$;
\item[(ii)] for every $Y$ the unique measure attaining the dual value is $\mu^\ast=(\tfr12,\tfr12)$ on the two worst-shock corner rows, with margin $-\varepsilon^\ast(Y)$ when $Y>\tfr{27}5$; the four-atom bundle (both shocks) has the same value and its optimal measure charges only the worst-shock atoms;
\item[(iii)] the Farkas pair of Theorem~\ref{thm:polyhedral} is $\lambda=(\tfr12,\tfr12)$, $\nu=\tfr12$ on the demand facet $-u_1-u_2\le-2$ and $\nu=0$ on the box facets: $\sum\lambda_av_a=(-\tfr12,-\tfr12)=\tfr12\,(-1,-1)$ and $\sum\lambda_ac_a+b^\top\nu=\tfr{\capi{1}+\capi{2}}{2}-1=\varepsilon^\ast(Y)$;
\item[(iv)] on the band $Y\in(\tfr{27}5,\tfr{69}5)$ each corner is first-order viable alone ($\capi{i}(Y)\ge0$) and the pair is not: the obstruction is epistemic; it is institutional relative to the box $U_0=[0,2]^2$, on which the fibre is viable with margin $\min_i\capi{i}(Y)>0$.
\end{enumerate}
At $Y=6$: $\capi{1}=\tfr{11}{10}$, $\capi{2}=\tfr{39}{50}$, cap sum $\tfr{47}{25}<2$, margin $-\tfr3{50}$; at $Y=5$ the common action $(\tfr65,\tfr45)$ is safe and the maximiser $(\tfr{29}{25},\tfr{21}{25})$ has margin $\tfr1{25}$; at $Y=4$ the caps are $\tfr{13}{10}$ and $\tfr{49}{50}$ with margin $\tfr7{50}$.
\end{proposition}

\begin{proof}
(i) For $u\in U_2$, $\min_i(\capi{i}-u_i)\le\tfr12\sum_i(\capi{i}-u_i)\le\tfr12(\capi{1}+\capi{2}-2)$, with equality at $u=(\capi{1}-\varepsilon,\capi{2}-\varepsilon)$, $\varepsilon=\tfr12(\capi{1}+\capi{2}-2)$, i.e. at $u^\ast=\tfr12(\capi{1}-\capi{2}+2,\ \capi{2}-\capi{1}+2)=(\tfr{29}{25},\tfr{21}{25})$, which does not depend on $Y$ and lies in $U_2$. Substituting \eqref{eq:caps} gives $\tfr{27-5Y}{50}$, which vanishes at $\tfr{27}5$. (ii) For $\lambda=(\lambda_1,\lambda_2)$ the pooled maximum over $U_2$ is $\lambda_1\capi{1}+\lambda_2\capi{2}-\min_{u\in U_2}(\lambda_1u_1+\lambda_2u_2)=\lambda_1\capi{1}+\lambda_2\capi{2}-2\min(\lambda_1,\lambda_2)$, which on either side of $\lambda_1=\tfr12$ is strictly monotone in $\lambda_1$ because $\capi{1}-\capi{2}+2>0$ and $\capi{2}-\capi{1}+2>0$; the minimum is at $\lambda=(\tfr12,\tfr12)$. The four-atom statement is a direct computation (check B2.3). (iii) is read off \eqref{eq:farkas} with $A$ the facet matrix of $U_2$ and $b$ its right-hand side; the box facets receive multiplier $0$. (iv) $\capi{2}(Y)\ge0$ iff $Y\le\tfr{69}5$ and $\capi{1}\ge\capi{2}$; at corner $i$ alone the action $u=(0,2)$ or $(2,0)$ is admissible and safe. On $U_0$ the common action $u=0$ is safe with margin $\min_i\capi{i}(Y)$.
\end{proof}

\begin{proposition}[demand-floor law]\label{prop:floor}
Let $Y^\ast(D)=\sup\{Y\ge4:\ \calR_{U_D}(B_Y)\neq\emptyset\}$ be the critical aggregate stock under the floor $D\in[0,4]$. Then
\begin{equation}\label{eq:floorlaw}
Y^\ast(D)=\min\Bigl(\tfrac{77}{5}-5D,\ \tfrac{69}{5}\Bigr)\quad\text{for }D\le\tfr{57}{25},\qquad\text{and no fibre } Y\ge4 \text{ is viable for } D>\tfr{57}{25}.
\end{equation}
For $D<\tfr8{25}$ the threshold $\tfr{69}5$ is a physical obstruction (the cap of corner $2$ vanishes); for $D\in(\tfr8{25},\tfr{57}{25}]$ the threshold $\tfr{77}5-5D$ is epistemic and institutional, with slope $-5=-\bigl(2\cdot\tfr1{20}\cdot2\bigr)^{-1}$: each unit of demand floor costs five units of aggregate stock. $D=2$ recovers $Y^\ast=\tfr{27}5$.
\end{proposition}

\begin{proof}
The fibre is viable iff $\capi{i}(Y)\ge0$ for both $i$ and $\capi{1}(Y)+\capi{2}(Y)\ge D$ (the same argument as in Proposition~\ref{prop:benchmark}(i) with $2$ replaced by $D$). The first condition is $Y\le\tfr{69}5$, the second $Y\le\tfr{77}5-5D$ since $\capi{1}+\capi{2}=\tfr{77}{25}-\tfr Y5$. The two lines cross at $D=\tfr8{25}$, and $\tfr{77}5-5D\ge4$ iff $D\le\tfr{57}{25}$. The classification follows from Proposition~\ref{prop:benchmark}(iv) with $2$ replaced by $D$. The script verifies the zero, the sign change on both sides and the regime boundaries on a rational grid of $D$ (checks B4.1--B4.4).
\end{proof}

\begin{figure}[t]
\centering
\includegraphics[width=\textwidth]{fig2_benchmark.pdf}
\caption{The two-stock benchmark (Propositions~\ref{prop:benchmark} and~\ref{prop:floor}). (a) Corner caps, their sum and the demand floor $D=2$ as functions of the aggregate stock; the shaded band $(\tfr{27}5,\tfr{69}5)$ is the epistemic--institutional obstruction. (b) The control set $U_2$ and the cap box at $Y=6$: the box $[0,\tfr{11}{10}]\times[0,\tfr{39}{50}]$ lies strictly below the demand facet; the pooled row with weights $(\tfr12,\tfr12)$ is the line $u_1+u_2=\tfr{47}{25}$, margin $-\tfr3{50}$. (c) The demand-floor law $Y^\ast(D)=\min(\tfr{77}5-5D,\tfr{69}5)$ with its physical plateau, its epistemic--institutional slope and the region without viable fibres.}
\label{fig:bench}
\end{figure}

\begin{remark}[what the aggregate observation costs]\label{rem:aggregate}
Under full observation of $(x_1,x_2)$ every fibre point with $Y\le\tfr{69}5$ is first-order viable, since the atom is a single state and each corner has a safe action. The aggregate observation collapses the fibre into one atom and loses the whole band $(\tfr{27}5,\tfr{69}5]$ of aggregate stocks: the \emph{value of resolving the aggregate} is a band of width $\tfr{42}5$ in $Y$, and the certificate says why---both corners must be charged, and the demand facet carries the multiplier. The three levers of Section~\ref{sec:taxonomy} appear with prices: resolve the atom (epistemic), lift or lower the floor at five units of stock per unit of demand (institutional), or accept that beyond $\tfr{69}5$ corner $2$ is lost under any observation and any floor (physical).
\end{remark}
